\begin{lemma}[CauchyCompletionMonadLaw idempotent metric-monad route]
\label{lem:cauchy-completion-monad-law-idempotent-metric-monad-route}
Every accepted $\CauchyCompletionMonadLawUp$ consumer that reads the unit, bind, multiplication, or associativity rows must first expose the completion idempotent route of \autoref{ch:concrete-instances-cauchy-completion-idempotent-namecert} and the metric completion monad route of \autoref{ch:concrete-instances-metric-completion-monad-namecert}. The idempotent sibling supplies the completion-after-completion stabilization row, while the metric completion monad sibling supplies the unit and bind law surface; the combined route exports no quotient-completion equality, selected representative, countable choice, host equality of completions, or ambient $\RealUp$ completeness principle.
\end{lemma}
\begin{proof}
Project the displayed law carrier from \autoref{def:cauchy-completion-monad-law-carrier}. Its unit, bind, multiplication, and associativity rows are read as completion-law data only after the carrier, stream, dyadic, regular rational, transport, replay, provenance, and local naming rows have been displayed.

Now read the completion idempotent sibling. The chapter \autoref{ch:concrete-instances-cauchy-completion-idempotent-namecert} supplies the stabilization surface for applying completion after completion, so a monad-law consumer cannot treat the multiplication row as an unrecorded host equality.

Finally read the metric completion monad sibling. The chapter \autoref{ch:concrete-instances-metric-completion-monad-namecert} supplies the unit and bind law surface consumed by the same route, and \autoref{thm:cauchy-completion-monad-law-namecert-obligations} keeps those rows inside the finite packet. Hence the sibling dependency lattice is substantive and finite, with none of the excluded quotient, selection, choice, host-equality, or ambient-completeness endpoints.
\end{proof}
